\hypertarget{chapter-04-page-081}{}
\subsection{定义与例子}
\label{sec:ch4-definition-examples}

我们所关心的奇异积分是如下形式的算子:
\begin{equation}
\label{eq:ch4-singular-integral-definition}
 Tf(x)=\lim_{\varepsilon\to0}\int_{|y|>\varepsilon}
 \frac{\Omega(y')}{|y|^n}f(x-y)\,dy,
\end{equation}
其中 $\Omega$ 定义在单位球面 $S^{n-1}$ 上, 可积且平均值为零, 并且 $y'=y/|y|$.

在这些假设下, 对 Schwartz 函数, \eqref{eq:ch4-singular-integral-definition} 有定义, 因为它是 $f$ 与缓增分布 $\operatorname{p.v.}\Omega(x')/|x|^n$ 的卷积, 后者定义为
\begin{equation}
\label{eq:ch4-principal-value-distribution}
\begin{aligned}
 \operatorname{p.v.}\frac{\Omega(x')}{|x|^n}(\varphi)
 &=\lim_{\varepsilon\to0}\int_{|x|>\varepsilon}
   \frac{\Omega(x')}{|x|^n}\varphi(x)\,dx\\
 &=\int_{|x|<1}\frac{\Omega(x')}{|x|^n}
   [\varphi(x)-\varphi(0)]\,dx
   +\int_{|x|>1}\frac{\Omega(x')}{|x|^n}\varphi(x)\,dx.
\end{aligned}
\end{equation}
第二个等式来自 $\Omega$ 的平均值为零. 由于 $\varphi\in\mathcal S$, 两个积分都收敛. 例如, 也可以假设 \eqref{eq:ch4-principal-value-distribution} 中的 $\varphi$ 或 \eqref{eq:ch4-singular-integral-definition} 中的 $f$ 是具有紧支集的 $\alpha$ 阶 Lipschitz 函数.

\begin{Proposition}
\label{prop:ch4-zero-mean-necessary}
极限 \eqref{eq:ch4-singular-integral-definition} 或 \eqref{eq:ch4-principal-value-distribution} 存在的一个必要条件是 $\Omega$ 在 $S^{n-1}$ 上的平均值为零.
\end{Proposition}

\begin{Proof}
取 $f\in\mathcal S(\mathbb R^n)$, 使得当 $|x|<2$ 时 $f(x)=1$. 当 $|x|<1$ 时,
\[
 Tf(x)=\int_{|y|>1}\frac{\Omega(y')}{|y|^n}f(x-y)\,dy
 +\lim_{\varepsilon\to0}\int_{\varepsilon<|y|<1}
 \frac{\Omega(y')}{|y|^n}\,dy.
\]
第一个积分总是收敛, 而第二项等于
\[
 \lim_{\varepsilon\to0}
 \left(\int_{S^{n-1}}\Omega(y')\,d\sigma(y')\right)
 \log(1/\varepsilon),
\]
它仅在 $\Omega$ 在 $S^{n-1}$ 上的积分为零时才有限.
\end{Proof}

\hypertarget{chapter-04-page-082}{}
当 $n=1$ 时, 单位球面退化为 $1$ 与 $-1$ 两个点, $\Omega$ 在这两个点必须取相反的值. 因此在实直线上, 任何形如 \eqref{eq:ch4-singular-integral-definition} 的算子都是 Hilbert 变换的常数倍.

在高维情形, 我们考察 A. P. Calderón 与 A. Zygmund 给出的两个例子. 首先, 设 $f(x_1,x_2)$ 是平面上的质量密度, 则它在上半空间 $\mathbb R_+^3$ 中的 Newton 势为
\[
 u(x_1,x_2,x_3)=\int_{\mathbb R^2}
 \frac{f(y_1,y_2)}{[(x_1-y_1)^2+(x_2-y_2)^2+x_3^2]^{1/2}}\,dy_1dy_2.
\]
引力场的强度由势的偏导数给出. $x_3$ 方向的分量等价于我们已经考察过的 Poisson 积分的常数倍. 另外两个分量彼此类似. 例如, 对第一个分量, 形式上有
\[
 \lim_{x_3\to0}\frac{\partial u}{\partial x_1}(x_1,x_2,x_3)
 =-\int_{\mathbb R^2}\frac{f(y_1,y_2)(x_1-y_1)}
 {[(x_1-y_1)^2+(x_2-y_2)^2]^{3/2}}\,dy_1dy_2.
\]
这个积分一般不收敛, 但若 $f$ 光滑, 它作为主值存在, 并且确实等于 $\partial u/\partial x_1$ 的极限. 它对应于 $\mathbb R^2$ 中形如 \eqref{eq:ch4-singular-integral-definition} 的奇异积分, 其中 $\Omega(x')=-x_1/|x|$; 在极坐标中, 这就是 $-\cos\theta$.

第二个例子来自平面质量分布 $f(x_1,x_2)$ 所对应的对数势, 即方程 $\Delta u=f$ 的解:
\[
 u(x_1,x_2)=\iint_{\mathbb R^2}f(y_1,y_2)
 \log([(x_1-y_1)^2+(x_2-y_2)^2]^{1/2})\,dy_1dy_2.
\]
例如当 $f$ 有紧支集时, 这个积分绝对收敛, 可以在积分号下求一阶偏导. 但对于二阶偏导不能这样做. 可以证明, $\partial^2u/\partial x_1\partial x_2$ 由形如 \eqref{eq:ch4-singular-integral-definition} 的算子给出, 其中 $\Omega(x')=2x_1x_2/|x|^2$. 关于这个算子的更多信息见第~\ref{sec:ch4-operator-algebra}~节.
